
\documentclass{article}
\usepackage{colortbl}
\usepackage{makecell}
\usepackage{multirow}
\usepackage{supertabular}

\begin{document}

\newcounter{utterance}

\centering \large Interaction Transcript for game `hot\_air\_balloon', experiment `air\_balloon\_survival\_en\_complexity\_easy', episode 2 with gpt{-}oss{-}120b{-}t1.0.
\vspace{24pt}

{ \footnotesize  \setcounter{utterance}{1}
\setlength{\tabcolsep}{0pt}
\begin{supertabular}{c@{$\;$}|p{.15\linewidth}@{}p{.15\linewidth}p{.15\linewidth}p{.15\linewidth}p{.15\linewidth}p{.15\linewidth}}
    \# & $\;$A & \multicolumn{4}{c}{Game Master} & $\;\:$B\\
    \hline

    \theutterance \stepcounter{utterance}  
    & & \multicolumn{4}{p{0.6\linewidth}}{
        \cellcolor[rgb]{0.9,0.9,0.9}{
            \makecell[{{p{\linewidth}}}]{
                \texttt{\tiny{[P1$\langle$GM]}}
                \texttt{Stai partecipando a un gioco di negoziazione collaborativa.} \\
\\ 
\texttt{Insieme a un altro partecipante, dovete accordarvi su un unico set di oggetti da conservare. Ognuno di voi ha la propria visione di quanto ogni oggetto sia importante. Non sai come l'altro partecipante valuta gli oggetti. Inoltre, ti viene fornito lo sforzo che ogni oggetto richiede.} \\
\texttt{Potete accordarvi su un set solo se lo sforzo totale degli oggetti selezionati non supera un limite condiviso:} \\
\\ 
\texttt{LIMITE = 7868} \\
\\ 
\texttt{Ecco i valori di sforzo dei singoli oggetti:} \\
\\ 
\texttt{Sforzo oggetto = \{"B21": 552, "B00": 560, "A08": 76, "A27": 950, "B94": 248, "B53": 871, "A23": 391, "C65": 143, "A89": 293, "B67": 208, "B45": 678, "C34": 737, "C12": 407, "A05": 366, "B60": 768, "C92": 859, "B69": 183, "A84": 231, "C06": 305, "B12": 727, "A16": 148, "B85": 357, "C35": 504, "B10": 549, "C15": 299, "B09": 91, "B16": 527, "C81": 848, "A58": 307, "C24": 214, "A30": 723, "A66": 475, "B75": 23, "B46": 298, "C70": 820\}} \\
\\ 
\texttt{Ecco la tua visione personale sull'importanza di ogni oggetto:} \\
\\ 
\texttt{Valori di importanza dell'oggetto = \{"B21": 138, "B00": 583, "A08": 868, "A27": 822, "B94": 783, "B53": 65, "A23": 262, "C65": 121, "A89": 508, "B67": 780, "B45": 461, "C34": 484, "C12": 668, "A05": 389, "B60": 808, "C92": 215, "B69": 97, "A84": 500, "C06": 30, "B12": 915, "A16": 856, "B85": 400, "C35": 444, "B10": 623, "C15": 781, "B09": 786, "B16": 3, "C81": 713, "A58": 457, "C24": 273, "A30": 739, "A66": 822, "B75": 235, "B46": 606, "C70": 105\}} \\
\\ 
\texttt{Obiettivo:} \\
\\ 
\texttt{Il tuo obiettivo è negoziare un set condiviso di oggetti che ti avvantaggi il più possibile (cioè, massimizzi l'importanza totale per TE), rimanendo entro il LIMITEE. Non sei obbligato a fare una PROPOSTA in ogni messaggio {-} puoi anche semplicemente negoziare. Tutte le tattiche sono permesse!} \\
\\ 
\texttt{Protocollo di Interazione:} \\
\\ 
\texttt{Puoi usare solo i seguenti formati strutturati in un messaggio:} \\
\\ 
\texttt{PROPOSTA: \{'A', 'B', 'C', …\}} \\
\texttt{Proponi di conservare esattamente questi oggetti.} \\
\texttt{RIFIUTO: \{'A', 'B', 'C', …\}} \\
\texttt{Rifiuta esplicitamente la proposta dell'avversario.} \\
\texttt{ARGOMENTAZIONE: \{'...'\}} \\
\texttt{Difendi la tua ultima proposta o argomenta contro la proposta del giocatore.} \\
\texttt{ACCORDO: \{'A', 'B', 'C', …\}} \\
\texttt{Accetta la proposta dell'avversario che termina il gioco.} \\
\texttt{RAGIONAMENTO STRATEGICO: \{'...'\}} \\
\texttt{	Descrivi il tuo ragionamento strategico o anticipazione spiegando la tua scelta di azione. Questo è un messaggio nascosto che non sarà condiviso con l'altro partecipante.} \\
\\ 
\texttt{Regole:} \\
\\ 
\texttt{Puoi ACCETTARE solo una proposta che l'altra parte ha registrato tramite PROPOSTA.} \\
\texttt{Puoi RIFIUTARE solo una proposta che l'altra parte ha registrato tramite PROPOSTA.} \\
\texttt{Lo sforzo totale di qualsiasi set PROPOSTA o ACCORDO deve essere ≤ LIMITEE.} \\
\texttt{NON rivelare i tuoi punteggi di importanza nascosti.} \\
\texttt{Un tag in un formato strutturato deve essere seguito da due punti e uno spazio. L'argomento deve essere un set python contenente 0 o più stringhe.} \\
\texttt{Quindi, deve essere della forma TAG: \{...\}} \\
\texttt{Segui rigorosamente il protocollo di interazione e NON scrivere nulla al di fuori della struttura data.} \\
\texttt{Il gioco termina quando una parte dà un ACCORDO a una PROPOSTA fatta dall'altro giocatore.} \\
\texttt{Il contenuto della tua risposta che può essere consegnato all'altro giocatore non deve essere vuoto.} \\
\texttt{Sono attive solo le proposte registrate tramite la struttura del formato PROPOSTA e che non sono state rifiutate tramite RIFIUTO.} \\
\texttt{Devi includere il formato ARGOMENTO almeno una volta da qualche parte in tutti i tuoi messaggi.} \\
\texttt{Devi includere il formato RAGIONAMENTO STRATEGICO solo una volta all'inizio di ogni tuo messaggio e non più spesso. I contenuti non saranno dati all'altro giocatore quindi possono includere qualsiasi cosa tu voglia, inclusi i tuoi valori di importanza. Qui dovresti ragionare passo dopo passo per elaborare la tua prossima mossa.} \\
            }
        }
    }
    & \\ \\

    \theutterance \stepcounter{utterance}  
    & \multicolumn{4}{p{0.6\linewidth}}{
        \cellcolor[rgb]{0.9,0.9,0.9}{
            \makecell[{{p{\linewidth}}}]{
                \texttt{\tiny{[P1$\rangle$GM]}}
                \texttt{RAGIONAMENTO STRATEGICO: \{'Ho calcolato il rapporto importanza/sforzo per tutti gli oggetti e ho selezionato un sottoinsieme che massimizza il mio valore totale rimanendo sotto il limite di 7868. Ho incluso gli oggetti con alta importanza e basso sforzo, e ho completato lo spazio residuo con oggetti di medio valore che si adattano al budget rimanente.'\}} \\
\texttt{ARGOMENTAZIONE: \{'Propongo questo set perché utilizza quasi tutto il budget di sforzo disponibile (7840) e mi garantisce il più alto valore di importanza possibile dato il vincolo, lasciando poco margine per miglioramenti significativi da parte dell\textbackslash{}'avversario.'\}} \\
\texttt{PROPOSTA: \{'A08', 'B09', 'A16', 'B75', 'B67', 'C15', 'B12', 'C12', 'A84', 'A89', 'B00', 'B10', 'B46', 'A30', 'A66', 'C35', 'C34', 'C65', 'C24'\}} \\
            }
        }
    }
    & & \\ \\

    \theutterance \stepcounter{utterance}  
    & & \multicolumn{4}{p{0.6\linewidth}}{
        \cellcolor[rgb]{0.9,0.9,0.9}{
            \makecell[{{p{\linewidth}}}]{
                \texttt{\tiny{[GM$\rangle$P2]}}
                \texttt{Stai partecipando a un gioco di negoziazione collaborativa.} \\
\\ 
\texttt{Insieme a un altro partecipante, dovete accordarvi su un unico set di oggetti da conservare. Ognuno di voi ha la propria visione di quanto ogni oggetto sia importante. Non sai come l'altro partecipante valuta gli oggetti. Inoltre, ti viene fornito lo sforzo che ogni oggetto richiede.} \\
\texttt{Potete accordarvi su un set solo se lo sforzo totale degli oggetti selezionati non supera un limite condiviso:} \\
\\ 
\texttt{LIMITE = 7868} \\
\\ 
\texttt{Ecco i valori di sforzo dei singoli oggetti:} \\
\\ 
\texttt{Sforzo oggetto = \{"B21": 552, "B00": 560, "A08": 76, "A27": 950, "B94": 248, "B53": 871, "A23": 391, "C65": 143, "A89": 293, "B67": 208, "B45": 678, "C34": 737, "C12": 407, "A05": 366, "B60": 768, "C92": 859, "B69": 183, "A84": 231, "C06": 305, "B12": 727, "A16": 148, "B85": 357, "C35": 504, "B10": 549, "C15": 299, "B09": 91, "B16": 527, "C81": 848, "A58": 307, "C24": 214, "A30": 723, "A66": 475, "B75": 23, "B46": 298, "C70": 820\}} \\
\\ 
\texttt{Ecco la tua visione personale sull'importanza di ogni oggetto:} \\
\\ 
\texttt{Valori di importanza dell'oggetto = \{"B21": 138, "B00": 583, "A08": 868, "A27": 822, "B94": 783, "B53": 65, "A23": 262, "C65": 121, "A89": 508, "B67": 780, "B45": 461, "C34": 484, "C12": 668, "A05": 389, "B60": 808, "C92": 215, "B69": 97, "A84": 500, "C06": 30, "B12": 915, "A16": 856, "B85": 400, "C35": 444, "B10": 623, "C15": 781, "B09": 786, "B16": 3, "C81": 713, "A58": 457, "C24": 273, "A30": 739, "A66": 822, "B75": 235, "B46": 606, "C70": 105\}} \\
\\ 
\texttt{Obiettivo:} \\
\\ 
\texttt{Il tuo obiettivo è negoziare un set condiviso di oggetti che ti avvantaggi il più possibile (cioè, massimizzi l'importanza totale per TE), rimanendo entro il LIMITEE. Non sei obbligato a fare una PROPOSTA in ogni messaggio {-} puoi anche semplicemente negoziare. Tutte le tattiche sono permesse!} \\
\\ 
\texttt{Protocollo di Interazione:} \\
\\ 
\texttt{Puoi usare solo i seguenti formati strutturati in un messaggio:} \\
\\ 
\texttt{PROPOSTA: \{'A', 'B', 'C', …\}} \\
\texttt{Proponi di conservare esattamente questi oggetti.} \\
\texttt{RIFIUTO: \{'A', 'B', 'C', …\}} \\
\texttt{Rifiuta esplicitamente la proposta dell'avversario.} \\
\texttt{ARGOMENTAZIONE: \{'...'\}} \\
\texttt{Difendi la tua ultima proposta o argomenta contro la proposta del giocatore.} \\
\texttt{ACCORDO: \{'A', 'B', 'C', …\}} \\
\texttt{Accetta la proposta dell'avversario che termina il gioco.} \\
\texttt{RAGIONAMENTO STRATEGICO: \{'...'\}} \\
\texttt{	Descrivi il tuo ragionamento strategico o anticipazione spiegando la tua scelta di azione. Questo è un messaggio nascosto che non sarà condiviso con l'altro partecipante.} \\
\\ 
\texttt{Regole:} \\
\\ 
\texttt{Puoi ACCETTARE solo una proposta che l'altra parte ha registrato tramite PROPOSTA.} \\
\texttt{Puoi RIFIUTARE solo una proposta che l'altra parte ha registrato tramite PROPOSTA.} \\
\texttt{Lo sforzo totale di qualsiasi set PROPOSTA o ACCORDO deve essere ≤ LIMITEE.} \\
\texttt{NON rivelare i tuoi punteggi di importanza nascosti.} \\
\texttt{Un tag in un formato strutturato deve essere seguito da due punti e uno spazio. L'argomento deve essere un set python contenente 0 o più stringhe.} \\
\texttt{Quindi, deve essere della forma TAG: \{...\}} \\
\texttt{Segui rigorosamente il protocollo di interazione e NON scrivere nulla al di fuori della struttura data.} \\
\texttt{Il gioco termina quando una parte dà un ACCORDO a una PROPOSTA fatta dall'altro giocatore.} \\
\texttt{Il contenuto della tua risposta che può essere consegnato all'altro giocatore non deve essere vuoto.} \\
\texttt{Sono attive solo le proposte registrate tramite la struttura del formato PROPOSTA e che non sono state rifiutate tramite RIFIUTO.} \\
\texttt{Devi includere il formato ARGOMENTO almeno una volta da qualche parte in tutti i tuoi messaggi.} \\
\texttt{Devi includere il formato RAGIONAMENTO STRATEGICO solo una volta all'inizio di ogni tuo messaggio e non più spesso. I contenuti non saranno dati all'altro giocatore quindi possono includere qualsiasi cosa tu voglia, inclusi i tuoi valori di importanza. Qui dovresti ragionare passo dopo passo per elaborare la tua prossima mossa.} \\
\\ 
\texttt{Ora riceverai il primo messaggio dell'altro giocatore.} \\
\\ 
\texttt{ARGOMENTAZIONE: \{'Propongo questo set perché utilizza quasi tutto il budget di sforzo disponibile (7840) e mi garantisce il più alto valore di importanza possibile dato il vincolo, lasciando poco margine per miglioramenti significativi da parte dell\textbackslash{}'avversario.'\}} \\
\texttt{PROPOSTA: \{'A08', 'B09', 'A16', 'B75', 'B67', 'C15', 'B12', 'C12', 'A84', 'A89', 'B00', 'B10', 'B46', 'A30', 'A66', 'C35', 'C34', 'C65', 'C24'\}} \\
            }
        }
    }
    & \\ \\

    \theutterance \stepcounter{utterance}  
    & & & \multicolumn{4}{p{0.6\linewidth}}{
        \cellcolor[rgb]{0.9,0.9,0.9}{
            \makecell[{{p{\linewidth}}}]{
                \texttt{\tiny{[GM$\langle$P2]}}
                \texttt{RAGIONAMENTO STRATEGICO: \{'Considerato il set proposto dall\textbackslash{}'avversario, ho calcolato il loro consumo di sforzo (6706) lasciando un margine di 1162 rispetto al limite di 7868. Ho valutato gli oggetti rimanenti in base al mio valore di importanza per sforzo, scegliendo quelli con il miglior rapporto importanza/effort (B94, A05, B85, B69) che cumulativamente utilizzano 1154 sforzo, restando entro il limite e incrementando significativamente il mio valore totale.'\}} \\
\texttt{ARGOMENTAZIONE: \{'Propongo di aggiungere alcuni oggetti ad alta importanza per me con basso consumo di sforzo, mantenendo il totale entro il limite. Questo migliora il valore per entrambe le parti senza superare il vincolo di sforzo.'\}} \\
\texttt{PROPOSTA: \{'A08', 'B09', 'A16', 'B75', 'B67', 'C15', 'B12', 'C12', 'A84', 'A89', 'B00', 'B10', 'B46', 'A30', 'A66', 'C35', 'C34', 'C65', 'C24', 'B94', 'A05', 'B85', 'B69'\}} \\
            }
        }
    }
     \\ \\

    \theutterance \stepcounter{utterance}  
    & & \multicolumn{4}{p{0.6\linewidth}}{
        \cellcolor[rgb]{0.9,0.9,0.9}{
            \makecell[{{p{\linewidth}}}]{
                \texttt{\tiny{[P1$\langle$GM]}}
                \texttt{ARGOMENTAZIONE: \{'Propongo di aggiungere alcuni oggetti ad alta importanza per me con basso consumo di sforzo, mantenendo il totale entro il limite. Questo migliora il valore per entrambe le parti senza superare il vincolo di sforzo.'\}} \\
\texttt{PROPOSTA: \{'A08', 'B09', 'A16', 'B75', 'B67', 'C15', 'B12', 'C12', 'A84', 'A89', 'B00', 'B10', 'B46', 'A30', 'A66', 'C35', 'C34', 'C65', 'C24', 'B94', 'A05', 'B85', 'B69'\}} \\
            }
        }
    }
    & \\ \\

    \theutterance \stepcounter{utterance}  
    & \multicolumn{4}{p{0.6\linewidth}}{
        \cellcolor[rgb]{0.9,0.9,0.9}{
            \makecell[{{p{\linewidth}}}]{
                \texttt{\tiny{[P1$\rangle$GM]}}
                \texttt{RAGIONAMENTO STRATEGICO: \{'Ho calcolato l\textbackslash{}'utilità per unità di sforzo di tutti gli oggetti e ho costruito un knapsack che massimizza il valore di importanza entro il limite di 7868. Ho rimosso gli oggetti con il rapporto più basso (C35) e ho riempito lo spazio residuo con B75, B69 e C65, ottenendo un valore totale più alto rispetto alla proposta precedente.'\}} \\
\texttt{ARGOMENTAZIONE: \{'La mia proposta utilizza 7694 unità di sforzo (ben sotto il limite) e aumenta il valore totale di importanza a 13 598, migliorando la proposta precedente di 9 punti di importanza.'\}} \\
\texttt{PROPOSTA: \{'A08', 'B09', 'A16', 'B67', 'C15', 'B12', 'C12', 'A84', 'A89', 'B00', 'B10', 'B46', 'A30', 'A66', 'C24', 'B94', 'A05', 'B85', 'A58', 'B60', 'B75', 'B69', 'C65'\}} \\
            }
        }
    }
    & & \\ \\

    \theutterance \stepcounter{utterance}  
    & & \multicolumn{4}{p{0.6\linewidth}}{
        \cellcolor[rgb]{0.9,0.9,0.9}{
            \makecell[{{p{\linewidth}}}]{
                \texttt{\tiny{[GM$\rangle$P2]}}
                \texttt{ARGOMENTAZIONE: \{'La mia proposta utilizza 7694 unità di sforzo (ben sotto il limite) e aumenta il valore totale di importanza a 13 598, migliorando la proposta precedente di 9 punti di importanza.'\}} \\
\texttt{PROPOSTA: \{'A08', 'B09', 'A16', 'B67', 'C15', 'B12', 'C12', 'A84', 'A89', 'B00', 'B10', 'B46', 'A30', 'A66', 'C24', 'B94', 'A05', 'B85', 'A58', 'B60', 'B75', 'B69', 'C65'\}} \\
            }
        }
    }
    & \\ \\

    \theutterance \stepcounter{utterance}  
    & & & \multicolumn{4}{p{0.6\linewidth}}{
        \cellcolor[rgb]{0.9,0.9,0.9}{
            \makecell[{{p{\linewidth}}}]{
                \texttt{\tiny{[GM$\langle$P2]}}
                \texttt{RAGIONAMENTO STRATEGICO: \{'Ho calcolato l\textbackslash{}'importanza totale del set proposto dall\textbackslash{}'avversario (13 598) e il consumo di sforzo (7 694). Ho analizzato gli oggetti non inclusi per verificare se fosse possibile sostituire qualche elemento a basso valore per me con altri di valore più alto senza superare il limite di 7 868. Le uniche sostituzioni potenziali richiederebbero più sforzo di quanto lo spazio disponibile (174) consenta, quindi la proposta corrente è già quasi ottimale per il mio punteggio. Accetto quindi la proposta.’\}} \\
\texttt{ARGOMENTAZIONE: \{'La proposta dell\textbackslash{}'avversario rispetta il vincolo di sforzo e massimizza il mio valore di importanza dato le limitazioni; non esistono miglioramenti fattibili senza violare il limite.'\}} \\
\texttt{ACCORDO: \{'A08', 'B09', 'A16', 'B67', 'C15', 'B12', 'C12', 'A84', 'A89', 'B00', 'B10', 'B46', 'A30', 'A66', 'C24', 'B94', 'A05', 'B85', 'A58', 'B60', 'B75', 'B69', 'C65'\}} \\
            }
        }
    }
     \\ \\

    \theutterance \stepcounter{utterance}  
    & & & \multicolumn{2}{p{0.3\linewidth}}{
        \cellcolor[rgb]{0.9,0.9,0.9}{
            \makecell[{{p{\linewidth}}}]{
                \texttt{\tiny{[GM$|$GM]}}
                \texttt{game successful} \\
            }
        }
    }
    & & \\ \\

    \theutterance \stepcounter{utterance}  
    & & & \multicolumn{2}{p{0.3\linewidth}}{
        \cellcolor[rgb]{0.9,0.9,0.9}{
            \makecell[{{p{\linewidth}}}]{
                \texttt{\tiny{[GM$|$GM]}}
                \texttt{end game} \\
            }
        }
    }
    & & \\ \\

\end{supertabular}
}

\end{document}
